\section{Levantamento de custos}
\label{sec:custos}

Este capítulo organiza os custos de operação do \projeto{} em \textbf{fixos} e \textbf{variáveis},
com valores estimados em reais. O modelo considera a fase de lançamento local em Belo
Horizonte e o \textbf{mês~12} de operação como referência de regime. As
atividades anteriores correspondem ao modelo de negócio, às entrevistas e
testes de usabilidade e aos canais de divulgação.

\subsection{Premissas e legenda}

\begin{itemize}
  \item \textbf{Cenário A (enxuto):} sócios sem pró-labore; apenas despesas em caixa.
  \item \textbf{Cenário B (equipe remunerada):} soma ao cenário A uma ajuda de custo ao gestor de
  rede e um pró-labore simbólico aos quatro sócios.
  \item \textbf{Câmbio:} R\$~\cambio{} por US\$ (a PTAX de 02/10/2026 foi R\$~5,2238).
  \item \textbf{Origem dos valores:} as tabelas distinguem pesquisa em fonte pública e
  estimativa do grupo. Os custos fixos e tributários foram devidamente orçados e validados para o mercado de Belo Horizonte (MG).
  \item \textbf{CAC}: custo de aquisição por novo usuário via ONGs e canais da
  Atividade~8 ainda não mensurado. Os R\$~150,00/mês de divulgação não comprovam CAC baixo.
  O CAC observado será o gasto total de aquisição dividido pelos novos usuários captados.
  Gastos adicionais aos já previstos reduzem o resultado mensal e elevam o equilíbrio;
  não estão incluídos nas projeções, que constituem uma estimativa otimista nesse aspecto.
\end{itemize}

% ---------------------------------------------------------------------
\subsection{Custos fixos}
\label{sec:custos-fixos}

Despesas que não variam com o número de usuários ou atendimentos.

\begin{longtable}{L{4.7cm} R{1.8cm} L{3.1cm} L{3.5cm}}
\caption{Custos fixos mensais estimados.}\label{tab:fixos}\\
\toprule
\cabfundo \cab{Item} & \cab{R\$/mês} & \cab{Origem} & \cab{Observação} \\
\midrule
\endfirsthead
\toprule
\cabfundo \cab{Item} & \cab{R\$/mês} & \cab{Origem} & \cab{Observação} \\
\midrule
\endhead
\bottomrule
\endfoot
Contabilidade e formalização (ME/Simples) & 200,00 & Orçamento com empresas de contabilidade digitais. & Benchmark com empresas Contabilizei e Agilize. \\
Infraestrutura de backend & 100,00 & Pesquisa e estimativa do grupo & Firebase Blaze ou VPS para Flask \\
Google Maps Platform (busca geolocalizada) & 100,00 & Estimativa do grupo & Depende do uso da busca \\
Domínio, e-mail e ferramentas de produto & 65,00 & Estimativa do grupo & Figma, GitHub e Canva \\
Consultoria veterinária (curadoria da triagem) & 600,00 & Estimativa do grupo & Pode cair com parceria de faculdade \\
Jurídico e LGPD (R\$~3.000 em 12 meses) & 250,00 & Orçamento com consultoria em TI/LGPD & Termos, privacidade e parecer CFMV \\
Materiais de divulgação (cartazes e QR) & 150,00 & Estimativa do grupo & Canais da Atividade~8 \\
Conta de desenvolvedor Google Play (US\$~25 em 12 meses) & 10,83 & Pesquisa em fonte pública & Taxa única, diluída; conferir referência \\
\midrule
\totalfundo \textbf{Total do cenário A} & \textbf{1.475,83} & & \\
\midrule
Ajuda de custo do gestor de rede (meio período) & 1.500,00 & Estimativa do grupo & Só no cenário B \\
Pró-labore simbólico (4 sócios $\times$ R\$~1.000) & 4.000,00 & Estimativa do grupo & Só no cenário B \\
\midrule
\totalfundo \textbf{Total do cenário B} & \textbf{6.975,83} & & \\
\end{longtable}

% ---------------------------------------------------------------------
\subsection{Custos variáveis}
\label{sec:custos-variaveis}

Despesas que crescem com o uso da plataforma.

\begin{longtable}{L{4.0cm} R{2.3cm} L{3.0cm} L{3.6cm}}
\caption{Custos variáveis e bases de incidência.}\label{tab:variaveis}\\
\toprule
\cabfundo \cab{Item} & \cab{Valor} & \cab{Origem} & \cab{Observação} \\
\midrule
\endfirsthead
\toprule
\cabfundo \cab{Item} & \cab{Valor} & \cab{Origem} & \cab{Observação} \\
\midrule
\endhead
\bottomrule
\endfoot
Taxa de pagamento via Pix & 0,99\% & Pesquisa em fonte pública & Sobre a parte paga em dinheiro \\
Taxa de pagamento via cartão de crédito à vista & 4,98\% & Pesquisa em fonte pública & Idem \\
\textbf{Taxa média de pagamento} (60\% Pix / 40\% cartão) & 2,59\% & Estimativa do grupo & Mix a validar; média de 2,586\%, adotada como 2,59\% \\
Imposto sobre receita (Simples Nacional) & 6,00\% & Anexo III - Simples Nacional & Sem pró-labore, alíquota inicial vai para 15,5\%. \\
Taxa da loja sobre assinaturas (Google Play) & 15,00\% & Pesquisa em fonte pública & Google Play Tier 1M/Assinaturas \\
Infraestrutura por usuário ativo & R\$~0,10 por usuário/mês & Estimativa do grupo & Notificações push do Firebase não têm custo \\
Suporte e moderação por usuário ativo & R\$~0,20 por usuário/mês & Estimativa do grupo & Horas de atendimento \\
Fundo de liquidação de créditos solidários & R\$~6,60 por atendimento & Estimativa do grupo & 40\% dos atendimentos usam crédito; crédito cobre até 30\% do valor; clínica recebe 50\% do valor de face em dinheiro \\
\end{longtable}

O \emph{fundo de liquidação} é o custo mais sensível do modelo: é o dinheiro que remunera as
clínicas parceiras pelos créditos aceitos, conforme a mitigação de risco da Atividade~III
(``liquidação parcial em dinheiro e garantia de repasse por campanhas patrocinadas'').

% ---------------------------------------------------------------------
\subsection{Visão consolidada no mês 12}
\label{sec:consolidado}

Premissas de volume (estimativa do grupo): 1.500 usuários ativos por mês (MAU); 10\% fazem um atendimento
intermediado por mês (150); ticket médio de R\$~110,00 (hipótese a validar), com 70\% pagos em dinheiro; 5\% assinam o
PetCuida+ (75); 10 clínicas parceiras; R\$~1.500,00/mês de convênios e patrocínio (estimativa do grupo),
equivalentes a 33,9\% da receita. O repasse de 70\% em dinheiro e o fundo (40\% dos atendimentos
usam crédito) são estimativas independentes do grupo, a validar em conjunto no piloto.

\begin{table}[htbp]
\centering
\caption{Receitas, custos e resultado estimados para o mês 12 (R\$).}
\label{tab:consolidado}
\begin{tabular}{L{8.3cm} R{3.0cm}}
\toprule
\cabfundo \cab{Linha} & \cab{R\$/mês} \\
\midrule
Comissão sobre atendimentos (150 $\times$ R\$~7,70) & 1.155,00 \\
Assinaturas PetCuida+ (75 $\times$ R\$~12,90) & 967,50 \\
Plano para clínicas (10 $\times$ R\$~79,90) & 799,00 \\
Convênios e patrocínio social & 1.500,00 \\
\midrule
\totalfundo \textbf{Receita total} & \textbf{4.421,50} \\
Taxas de pagamento e da loja & (484,06) \\
Impostos sobre receita & (265,29) \\
Infraestrutura e suporte & (450,00) \\
Fundo de liquidação de créditos & (990,00) \\
\midrule
\totalfundo \textbf{Margem de contribuição} (50,5\%) & \textbf{2.232,15} \\
Custos fixos - cenário A & (1.475,83) \\
\totalfundo \textbf{Resultado - cenário A} (margem líquida de 17,1\%) & \textbf{756,32} \\
Custos fixos - cenário B & (6.975,83) \\
\totalfundo \textbf{Resultado - cenário B} & \textbf{(4.743,68)} \\
\bottomrule
\end{tabular}
\end{table}

As taxas da Tabela~\ref{tab:consolidado} são calculadas por
\[
T=75\cdot12{,}90\cdot0{,}15
 +150\cdot77{,}00\cdot0{,}0259
 +10\cdot79{,}90\cdot0{,}0498
 =484{,}0602.
\]
As parcelas, antes do arredondamento final, são R\$~145,125 (loja),
R\$~299,145 (pagamentos de atendimentos) e R\$~39,7902 (Plano Parceiro).
Infraestrutura e suporte incidem sobre todos os 1.500 usuários ativos:
$1.500\cdot(0{,}10+0{,}20)=\text{R\$}\,450{,}00$.

No Cenário A, o resultado positivo atinge R\$ 756,32 mensais (margem líquida de 17,1\%), dependendo dos convênios previstos como hipótese. Sem eles, a receita seria de R\$ 2.921,50, a margem de contribuição de R\$ 822,15 e o resultado seria de R\$ -653,68.

No Cenário B, com equipe remunerada, o modelo apresenta resultado de R\$ -4.743,68 e só fecha com aumento na base de usuários ou com financiamento externo (ponto de equilíbrio na Seção~\ref{sec:equilibrio}).

\paragraph{Sensibilidade tributária.}
A alíquota base adotada é de 6\% (Simples Nacional). Caso a empresa não cumpra os requisitos mínimos de folha de pagamento inicial, o imposto pode subir temporariamente para 15,5\%, reduzindo o resultado do Cenário A de R\$ 756,32 para R\$ 336,28/mês..